\DocumentMetadata{
  pdfversion=2.0,
  pdfstandard=ua-2,
  lang=en-US,
  tagging=on,
  tagging-setup={math/setup=mathml-SE,tabsorder=structure}
}
\documentclass[11pt,letterpaper]{article}
\usepackage[margin=0.68in,includefoot]{geometry}
\usepackage{newcomputermodern}
\usepackage{amsmath,amscd,mathtools}
\providecommand{\square}{\mdlgwhtsquare}
\usepackage[hidelinks]{hyperref}
\hypersetup{
  pdftitle={Numerical Analysis PhD Exam, August 2017, Part 2},
  pdfauthor={University of Florida Department of Mathematics},
  pdfsubject={Graduate mathematics examination},
  pdfdisplaydoctitle=true,
  bookmarksopen=true
}
\setcounter{secnumdepth}{0}
\setlength{\parindent}{0pt}
\setlength{\parskip}{0.28em}
\setlength{\emergencystretch}{3em}
\setlength{\leftmargini}{2.15em}
\setlength{\leftmarginii}{2.65em}
\setlength{\leftmarginiii}{3em}
\renewcommand{\labelenumi}{\textbf{\theenumi.}}
\pagestyle{empty}
\raggedbottom
\allowdisplaybreaks
\newenvironment{examproblems}[1]{%
  \begin{enumerate}%
  \setcounter{enumi}{#1}%
  \setlength{\topsep}{0.35em}%
  \setlength{\partopsep}{0pt}%
  \setlength{\itemsep}{0.48em}%
  \setlength{\parsep}{0.18em}%
}{\end{enumerate}}
\newenvironment{examsubparts}{%
  \begin{enumerate}%
  \setlength{\topsep}{0.18em}%
  \setlength{\partopsep}{0pt}%
  \setlength{\itemsep}{0.16em}%
  \setlength{\parsep}{0.1em}%
}{\end{enumerate}}
\newenvironment{examinstructions}{%
  \par\begingroup\itshape
}{%
  \par\endgroup\vspace{0.18em}
}
\makeatletter
\renewcommand\section{\@startsection{section}{1}{0pt}%
  {-1.25ex plus -.25ex minus -.1ex}%
  {0.55ex plus .1ex}%
  {\normalfont\large\bfseries}}
\renewcommand{\@maketitle}{%
  \newpage\null\vskip -1.2em%
  {\footnotesize\bfseries
    \href{https://www.ufl.edu/}{University of Florida}, \href{https://math.ufl.edu/}{Department of Mathematics}\par}%
  \vskip 0.18em%
  \tagstructbegin{tag=Title}%
    {\LARGE\bfseries\href{https://gma.math.ufl.edu/exams/numerical-analysis-phd/}{Numerical Analysis PhD Exam}, August 2017, Part 2\par}%
  \tagstructend%
  \vskip 0.35em%
  \tagmcbegin{artifact}%
    \hrule height 1pt%
  \tagmcend%
  \vskip 0.55em%
}
\makeatother
\title{Numerical Analysis PhD Exam, August 2017, Part 2}
\author{}
\date{}
\begin{document}
\maketitle
\thispagestyle{empty}
\begin{examinstructions}
Do all 5 (five) problems.
\end{examinstructions}
\begin{examproblems}{0}
\item
Let \(f(x)=e^x\).
\begin{examsubparts}
\item[\textnormal{(a)}]
Find the linear Taylor polynomial \(T_1(x)\) of \(f(x)\) expanded about \(x_0=1/2\) and give an estimate for the maximum of \(|T_1(x)-f(x)|\) on the interval \([0,1]\).
\item[\textnormal{(b)}]
Find the linear minimax approximation \(P_1(x)\) to \(f(x)\) on \([0,1]\) and find the maximum error of \(|P_1(x)-f(x)|\) on \([0,1]\).
\item[\textnormal{(c)}]
Find the linear least squares approximation to~\(f\) on \([0,1]\).
\end{examsubparts}
\item
The midpoint method for DE’s is 
\[
w_0=\alpha
\]
 
\[
w_{n+1}=w_n+hf\left(t_n+\frac{h}{2},w_n+\frac{h}{2}f(t_n,w_n)\right).
\]
 Apply this method to the IVP, \(y'=\lambda y\); \(y(0)=1\), with \(\lambda<0\) and find the constant~\(c\) so that \(|h\lambda|<c\) implies that \(w_n\to0\) as \(n\to\infty\).
\item
Derive this three-point formula for the second derivative. 
\[
f''(x_0)=\frac{1}{h^2}\bigl(f(x_0-h)-2f(x_0)+f(x_0+h)\bigr)-\frac{h^2}{12}f^{(4)}(\eta)
\]
 for some \(\eta\in[x_0-h,x_0+h]\).
\item
Let \(f(x)=2^x\). Let \(x_0=-1,x_1=0,x_2=1\). Find the total error bound for the degree two interpolating polynomial \(p_2(x)\) with these nodes and~\(f\) on the interval \([-1,1]\), i.e. derive a~\(K\) with 
\[
\max_{t\in[-1,1]}|f(t)-p_2(t)|<K.
\]
\item
Let \(g\in C^2([a,b])\) and \(p\in(a,b)\) with \(g(p)=p\), \(g'(p)=0\), \(g''(p)\ne0\).
\begin{examsubparts}
\item[\textnormal{(a)}]
Show there is an \(\epsilon>0\) so that for all \(x\in[p-\epsilon,p+\epsilon]\), we have \(g^n(x)\to p\) as \(n\to\infty\).
\item[\textnormal{(b)}]
With~\(\epsilon\) as in part (a), show that for all \(x\in[p-\epsilon,p+\epsilon]\), we have \(|g(x)-p|\le M|x-p|^2\), where 
\[
M=\max\{|g''(x)|:|x-p|\le\epsilon\}/2.
\]
\end{examsubparts}
\end{examproblems}
\end{document}
